\chapter{Footprint: JVM, nativo, memoria, GC}
\label{cap:24-footprint}

Questo capitolo risponde alla quarta domanda di ricerca del
\cref{cap:01-introduzione}: qual \`e il costo reale di un control plane, e come
cambia al variare della catena di build. \`E il capitolo con la maggiore densit\`a
di risultati negativi e di ipotesi rigettate, ed \`e riportato con quelle
incluse.

\section{Dove va la memoria di una JVM Spring}

Il punto di partenza \`e una funzione, non il control plane: \code{word-stats}
Java su Spring Boot, a 190\,MiB di memoria residente.

\begin{center}
\small
\begin{adjustbox}{max width=\textwidth}
\begin{tabular}{@{}lrr@{}}
\toprule
\textbf{Area} & \textbf{Dimensione} & \textbf{Quota} \\
\midrule
\textbf{heap in uso} & \textbf{22{,}4\,MB} & \textbf{12\,\%} \\
Metaspace (9\,375 classi caricate) & 53{,}0\,MB & 28\,\% \\
Compressed class space & 6{,}6\,MB & 3\,\% \\
CodeHeap (codice compilato JIT) & 13{,}5\,MB & 7\,\% \\
23 stack di thread & $\sim$23\,MB & 12\,\% \\
\emph{resto} (buffer diretti, GC, nativo) & $\sim$71\,MB & 37\,\% \\
\bottomrule
\end{tabular}
\end{adjustbox}
\end{center}

\begin{quote}
Lo heap --- i \emph{dati} --- \`e il 12\% del totale. Il resto \`e il runtime che
descrive se stesso: metadati delle classi, codice compilato, e i profili
conservati per ricompilarlo. Spring \`e costoso in memoria perch\'e carica
migliaia di classi e perch\'e una JVM porta con s\'e un compilatore, non perch\'e
allochi molto.
\end{quote}

Questa osservazione ha una conseguenza pratica immediata: ridurre l'uso di
memoria di un'applicazione Spring ottimizzando le allocazioni agisce sul 12\% del
problema. Le leve efficaci sono altre --- caricare meno classi, compilare in
anticipo, ridurre gli stack.

\`E anche la giustificazione quantitativa dei flag JVM del
\cref{cap:20-build-packaging}: \code{-Xss256k} agisce sui 23\,MB degli stack. Sui 13{,}5\,MB della CodeHeap
agiva \code{TieredStopAtLevel=1}, che il default di oggi non contiene pi\`u:
la latenza che C1 faceva perdere sul percorso caldo superava il risparmio
(\cref{cap:20-build-packaging}).

\section{Compilare nativamente}

\subsection{Sulla funzione}

\begin{center}
\small
\begin{adjustbox}{max width=\textwidth}
\begin{tabular}{@{}lrr@{}}
\toprule
 & \textbf{JVM} & \textbf{nativo} \\
\midrule
residente a riposo & 133\,MiB & \textbf{25{,}6\,MiB} \\
avvio & 1{,}046\,s & \textbf{0{,}066\,s} \\
residente sotto carico & 160\,MiB & 56\,MiB \\
\bottomrule
\end{tabular}
\end{adjustbox}
\end{center}

Il fattore cinque a riposo e il fattore sedici sull'avvio sono quanto ci si
aspetta dalla compilazione anticipata: sparisce il compilatore, spariscono i
metadati delle classi caricate dinamicamente.

Il dato interessante \`e l'ultimo. Sotto carico, la funzione Spring compilata
nativamente (56\,MiB) e la funzione \emph{lite} nativa (51{,}6\,MiB) quasi
convergono, pur differendo di un fattore dieci a riposo (25{,}6 contro
2{,}7\,MiB). L'interpretazione: \emph{il working set a caldo \`e dominato
dall'allocazione per richiesta, non dal framework}.

\`E un risultato che ridimensiona un'intuizione comune. La scelta di un framework
leggero fa una differenza enorme sull'impronta a riposo --- che conta per una
piattaforma che tiene molte funzioni idle --- e una differenza modesta
sull'impronta sotto carico.

\subsection{Sul control plane: il primo risultato \`e sbagliato}

La prima misura sembr\`o mostrare che il control plane nativo dimezzava sia
memoria sia CPU. Aveva dimezzato il \emph{lavoro}: 415\,184 richieste contro
865\,692, con p99 a 1\,696\,ms.

\`E l'errore di misura pi\`u istruttivo del progetto: il consumo di risorse per
unit\`a di tempo non \`e una metrica di efficienza se il sistema non sta
consegnando lo stesso lavoro.

\section{Il garbage collector dell'immagine nativa}

\subsection{La misura in isolamento}

Stesso backend, stesso carico, unica differenza la build del control plane:

\begin{center}
\small
\begin{adjustbox}{max width=\textwidth}
\begin{tabular}{@{}lrrr@{}}
\toprule
 & \textbf{JVM} & \textbf{nativo (serial)} & \textbf{nativo (G1)} \\
\midrule
throughput & 8\,205 rps & 4\,983 rps & \textbf{8\,085 rps} \\
mediana & 5{,}98\,ms & \textbf{4{,}71\,ms} & 6{,}08\,ms \\
p95 & 16{,}06\,ms & 27{,}95\,ms & 16{,}51\,ms \\
\textbf{massimo} & 92\,ms & \textbf{1{,}77\,s} & 123\,ms \\
avvio & 1{,}585\,s & 0{,}192\,s & \textbf{0{,}072\,s} \\
\bottomrule
\end{tabular}
\end{adjustbox}
\end{center}

Due osservazioni.

La \textbf{mediana della build nativa \`e migliore di quella della JVM}. Fra una
pausa e l'altra il codice compilato in anticipo \`e veloce; la perdita \`e
interamente sulla coda.

Il \textbf{massimo della build seriale \`e 1{,}77 secondi}, venti volte il
massimo della JVM. \`E la firma di un sistema che si ferma, non di uno che \`e
lento.

Il costo di CPU per richiesta era 1{,}62\,ms contro 1{,}32 --- la penalizzazione
ordinaria della compilazione anticipata, il 23\%, ben lontana dallo spiegare un
dimezzamento. E il sistema usava il 265\% di un 1\,100\% disponibile mentre
consegnava met\`a del lavoro: \emph{nulla era saturo, qualcosa si fermava}.

\subsection{La causa}

Corretto un difetto nel binder delle metriche del collector --- \code{JvmGcMetrics}
richiede notifiche che SubstrateVM non emette, e riportava uno zero costante ---
la causa \`e risultata diretta:

\begin{center}
\small
\begin{adjustbox}{max width=\textwidth}
\begin{tabular}{@{}lrr@{}}
\toprule
 & \textbf{JVM} & \textbf{nativo (serial)} \\
\midrule
raccolte complete & 2, $\sim$50\,ms ciascuna & \textbf{25, 883\,ms ciascuna} \\
tempo totale di raccolta & 6{,}6\,s & \textbf{22{,}1\,s su 60} \\
quota del tempo di parete & \textbf{0{,}9\,\%} & \textbf{31{,}8\,\%} \\
\bottomrule
\end{tabular}
\end{adjustbox}
\end{center}

Quasi un terzo del tempo speso a raccogliere.

\subsection{Tre ipotesi rigettate}

Il progetto ha testato e scartato tre spiegazioni alternative, e le riporta:

\begin{itemize}[leftmargin=*]
  \item \textbf{Limitare lo heap.} \code{-Xmx512m} ha peggiorato molto la
  situazione: 1\,543 rps e timeout a 30\,s.
  \item \textbf{Promozione prematura.} Ingrandire la generazione giovane non ha
  cambiato nulla. \emph{Il primo tentativo di questo test non era valido}:
  \code{-XX:MaxNewSize} \`e limitato da
  \code{-XX:MaximumYoungGenerationSizePercent}, il cui default \`e 10, quindi il
  flag non faceva nulla.
  \item \textbf{Buffer heap di Netty.} Entrambe le build preferiscono buffer
  diretti, entrambe usano l'allocatore adattivo, a nessuna delle due \`e
  disponibile \code{Unsafe}. Identiche.
\end{itemize}

Il secondo caso \`e riportato con la propria invalidazione. \`E buona pratica
sperimentale e vale la pena notarla: un test che non fa ci\`o che crede di fare
produce un risultato ``nessun effetto'' indistinguibile da quello vero.

\subsection{Il meccanismo}

La pausa scala con la \emph{dimensione} dello heap, non con i dati vivi:
563\,ms per uno heap di 512\,MB contro un insieme vivo di decine di megabyte. \`E
il comportamento di un collector che percorre l'intero spazio --- che \`e ci\`o
che fa l'unico collector utilizzabile di GraalVM Community: \code{--gc} accetta
\code{serial} ed \code{epsilon}, e nient'altro.

G1, disponibile in Oracle GraalVM, elimina il problema: 98{,}5\% del throughput
della JVM, coda a 123\,ms, avvio in 0{,}07\,s. Poich\'e Oracle GraalVM ha licenza
GFTC anzich\'e GPL, la build resta opt-in tramite \code{NATIVE\_GC=G1}.

\begin{damisurare}
\emph{Perch\'e} una raccolta completa costi 883\,ms con un insieme vivo piccolo
non \`e stato confermato. La pausa traccia la dimensione dello heap anzich\'e i
dati vivi, ma il meccanismo non \`e stato stabilito.
\todomis{meccanismo della pausa del collector seriale di SubstrateVM}
\end{damisurare}

\section{Osservare il collector: un problema di strumento}

Sotto G1 l'immagine nativa non registra alcun \code{GarbageCollectorMXBean},
quindi n\'e Micrometer n\'e il binder a polling riportano nulla --- sulla build
che si comporta meglio. \code{--enable-monitoring=jfr,jvmstat,jmxserver} non li
riporta in vita; \`e stato provato e non funziona.

JFR funziona, per una via indiretta. I tipi di evento specifici del collector
(\code{jdk.G1GarbageCollection}, \code{jdk.GCHeapSummary}, la famiglia
\code{GCPhase*}) sono dichiarati nella registrazione ed emettono zero eventi. Le
pause sono registrate come \textbf{operazioni della VM}, ciascuna con la propria
durata. Da una registrazione di 90\,s sotto carico:

\begin{center}
\small
\begin{adjustbox}{max width=\textwidth}
\begin{tabular}{@{}lrrrr@{}}
\toprule
\textbf{Operazione} & \textbf{Conteggio} & \textbf{Totale} & \textbf{Media} &
\textbf{Massimo} \\
\midrule
\code{G1 wrapper} & 43 & 835{,}1\,ms & 19{,}4\,ms & 97{,}6\,ms \\
\code{Collect for allocation} & 31 & 651{,}1\,ms & 21{,}0\,ms & 94{,}5\,ms \\
\code{Try init concurrent mark} & 2 & 37{,}3\,ms & 18{,}6\,ms & 22{,}5\,ms \\
\midrule
\textbf{totale} & \textbf{76} & \textbf{1\,524\,ms su 90\,s} & &
\textbf{1{,}69\,\%} \\
\bottomrule
\end{tabular}
\end{adjustbox}
\end{center}

Il risultato conferma la sezione precedente per una via indipendente ---
1{,}69\% del tempo di parete contro il 31{,}8\% del collector seriale --- e
fornisce pi\`u di quanto gli MXBean potrebbero: una durata per singola pausa, e
quindi il massimo, che i contatori interrogati periodicamente non danno neppure
dove funzionano.

\begin{center}
\small
\begin{adjustbox}{max width=\textwidth}
\begin{tabular}{@{}lll@{}}
\toprule
\textbf{Build} & \textbf{Conteggi e tempo totale} & \textbf{Pausa massima} \\
\midrule
JVM & Micrometer, nativamente & s\`i, dalle notifiche \\
nativo, serial & il binder a polling & no \\
nativo, G1 & \textbf{solo JFR} & \textbf{s\`i, per pausa} \\
\bottomrule
\end{tabular}
\end{adjustbox}
\end{center}

La conclusione \`e precisa: \emph{G1 non \`e inosservabile, \`e inosservabile
attraverso Prometheus}. Il dato richiede una registrazione estratta e analizzata
fuori banda anzich\'e un gauge da raccogliere. Adeguato per un'indagine, non per
un allarme.

\`E anche un limite reale dell'apparato sperimentale di questo progetto, che
altrove poggia interamente su Prometheus.

\section{Il livello di ottimizzazione}

\begin{center}
\small
\begin{adjustbox}{max width=\textwidth}
\begin{tabular}{@{}lrr@{}}
\toprule
 & \textbf{\code{-Os}} & \textbf{\code{-O3}} \\
\midrule
throughput & 8\,398 rps & \textbf{9\,251 rps} $(+10{,}2\%)$ \\
p95 & 15{,}64\,ms & 14{,}79\,ms \\
massimo & 148{,}7\,ms & 127{,}2\,ms \\
avvio & 0{,}173\,s & \textbf{0{,}068\,s} \\
residente a riposo & 120{,}6\,MiB & \textbf{39{,}3\,MiB} \\
dimensione immagine & 195\,MB & 397\,MB \\
\bottomrule
\end{tabular}
\end{adjustbox}
\end{center}

A \code{-O3} la build nativa non si limita a eguagliare la JVM: la supera ---
9\,251 contro 8\,205 richieste al secondo, con un avvio ventitr\'e volte pi\`u
rapido.

\subsection{Perch\'e l'immagine cresce}

Dal rapporto di build: l'area di codice triplica, da 47{,}78\,MB a 166{,}59\,MB,
a partire da \emph{meno} unit\`a di compilazione --- 84\,694 contro 105\,256.

Quell'accoppiamento \`e la firma dell'inlining. I metodi vengono copiati nei
chiamanti anzich\'e chiamati, quindi le unit\`a spariscono mentre ogni
sopravvissuta porta copie di ci\`o che ha assorbito; lo srotolamento dei cicli
aggiunge il resto. Stesso programma, pi\`u codice macchina --- che \`e anche
perch\'e \`e pi\`u veloce: nessun costo di chiamata, e ogni copia specializzata
per il proprio sito.

\subsection{Dove cade il costo, e dove non cade il guadagno}

Due qualificazioni, entrambe misurate.

Con un tetto di heap fissato, entrambi i livelli hanno raggiunto lo stesso picco
di 744\,MiB residenti, perch\'e il codice \`e mappato dal file anzich\'e copiato
nello heap: i 200\,MB in pi\`u sono un costo per il registry e per il disco, non
per il nodo.

Il guadagno di throughput \`e \emph{condizionale}: sotto un limite di 1\,GB le
due build sono risultate identiche a circa 3\,490 richieste al secondo, perch\'e
l\`i il collo di bottiglia \`e la raccolta e non la qualit\`a del codice.
\code{-O3} compra velocit\`a solo dove la memoria non \`e gi\`a il vincolo.

\subsection{Un avvertimento su PGO}

Il log di build riporta \code{PGO: off} a \code{-Os} e
\code{PGO: ML-inferred} a \code{-O3}: Oracle GraalVM applica profili inferiti da
modello quando ottimizza per velocit\`a. Parte del guadagno \`e quindi gi\`a una
forma di guida per profilo. Un eventuale PGO esplicito va misurato contro
\code{-O3}, non contro \code{-Os}, o gli verr\`a accreditato ci\`o che
\code{-O3} ha gi\`a raccolto.

\section{I limiti di memoria cambiano la risposta}

Tutte le cifre precedenti sono senza limite di memoria del container, dove ogni
build cresce finch\'e il suo collector non vede ragione di fermarsi.

\begin{table}[htbp]
\centering
\caption{Throughput e latenza sotto limite di memoria del container.}
\label{tab:memlimit}
\small
\begin{adjustbox}{max width=\textwidth}
\begin{tabular}{@{}llrrrr@{}}
\toprule
\textbf{Limite} & \textbf{Build} & \textbf{Throughput} & \textbf{p95} &
\textbf{Max} & \textbf{Memoria usata} \\
\midrule
500\,MB & JVM & 1\,515 rps & 324\,ms & 1{,}99\,s & 455 / 500 \\
        & nativo serial & 1\,370 rps & 32\,ms & 30{,}0\,s & 489 / 500 \\
        & \textbf{nativo G1, default} & \textbf{520 rps} & 61\,ms & 30{,}0\,s &
        \textbf{145 / 500} \\
        & nativo G1, \code{-Xmx350m} & 1\,445 rps & 96\,ms & 30{,}0\,s & 378 / 500 \\
\midrule
1\,GB & JVM & 2\,450 rps & 46\,ms & 1{,}34\,s & 805 / 1024 \\
      & \textbf{nativo G1, \code{-Xmx700m}} & \textbf{3\,089 rps} & 44\,ms &
      \textbf{135\,ms} & 736 / 1024 \\
\midrule
2\,GB & JVM & 2\,712 rps & 47{,}5\,ms & 451\,ms & 727 / 2048 \\
      & nativo G1, \code{-Xmx1400m} & 3\,400 rps & 39{,}5\,ms & 151\,ms & 1\,201 / 2048 \\
\midrule
nessuno & JVM & 8\,205 rps & 16{,}1\,ms & 92\,ms & 1{,}8\,GiB \\
        & nativo G1 & 8\,085 rps & 16{,}5\,ms & 123\,ms & 2{,}7\,GiB \\
\bottomrule
\end{tabular}
\end{adjustbox}
\end{table}

\subsection{G1 sotto un limite piccolo \`e una trappola}

Il tetto di default di G1 \`e il 25\% della RAM visibile. Rispetta il cgroup ---
il che \`e corretto --- ma poi si confina a 125\,MB di un container da 500\,MB,
lascia inutilizzato il 71\% della memoria, e crolla a 520 rps. Fissare
esplicitamente \code{-Xmx} quasi lo triplica.

\begin{quote}
\textbf{Con un limite di memoria, una build nativa G1 deve ricevere uno heap
esplicito.}
\end{quote}

\subsection{A 500\,MB non ci sta nulla}

Tutte le build perdono da cinque a sei volte il throughput senza vincolo, e
quelle native iniziano ad andare in timeout. Il rimedio non \`e un collector
migliore ma meno concorrenza offerta.

\subsection{Il ginocchio}

Il ginocchio della curva \`e fra 1 e 2\,GB. A 1\,GB con G1 e
\code{-Xmx700m} si ottengono 3\,089 rps con una coda di 135\,ms e nessun
fallimento, in un terzo della memoria che l'esecuzione senza vincoli aveva
occupato.

\subsection{Una configurazione da correggere}

Il control plane distribuito via Docker Compose fissa
\code{-XX:MaxRAMPercentage=70} senza \code{mem\_limit}, il che su un host da
12{,}5\,GB autorizza uno heap da 8{,}77\,GB. Fino a che non venga fissato un
limite, ogni cifra di memoria proveniente da quel percorso descrive
\emph{appetito}, non fabbisogno.

\`E la conclusione generale del capitolo, ed \`e metodologica prima che tecnica:
\textbf{misure di memoria senza un limite descrivono quanto un runtime \`e
disposto a prendere, non quanto gli serve.}

\section{Sintesi}

\begin{table}[htbp]
\centering
\caption{Quando conviene ciascuna build.}
\label{tab:build-scelta}
\small
\begin{tabularx}{\textwidth}{@{}lXX@{}}
\toprule
\textbf{Build} & \textbf{Conviene quando} & \textbf{Da evitare quando} \\
\midrule
JVM & memoria abbondante; si vuole osservabilit\`a del collector via Prometheus &
il tempo di avvio conta \\
nativo, serial & si tollerano pause al secondo & sotto carico sostenuto: 32\% del
tempo speso a raccogliere \\
nativo, G1, \code{-O3} & throughput e avvio contano; si accetta la licenza
Oracle & sotto limiti di memoria stretti senza \code{-Xmx} esplicito \\
\bottomrule
\end{tabularx}
\end{table}
